\subsection{A Long Divion Algorithm in $k[x_1, \dots, x_n]$} 

In the univariate case $k[x]$: when $f$ is divided by $g$, one can explain the remainder $r$ with ideals: 
\begin{center}
    $r$ is the polynomial if smallest degree with $f-r \in \underbrace{\langle g \rangle}_{quotient g hodden here}$.  \\
\end{center}

In the multivariate case, an ideal is usually \underline{not} generated by  just one polynomial. So, we will need to divide by a list of polynomials. 
The other difference to the univariate case: dependence of on the order. 

Given $f,g_1, \ldots g_s \in k[x_1, \dots, x_n]$ we look for a representation 
\[f = g_1 \cdot  q_1 + \ldots + g_s \cdot q_s + r\]
where $q_1, \ldots, q_s$ are quotients and $r$ is the remainder. \\[0.5em]
\textbf{Issue:} sometimes one cannot divide monomials. 
\[x_1^2 \cdot x_2 \succ_{lex} x_2 \cdot x_2^2\] 
but $x_1^2 \cdot x_2$ is not divisible by $x_2 \cdot x_2^2$. \\[0.5em]
\textbf{Way out:} put everything that can't be divided into the remainder. 

\clearpage 

\begin{algorithm} [title = Multivariate Division Algorithm]
    \textbf{Assumption:} $f, g_1, \ldots, g_s \in k[x_1, \dots, x_n], g_i \neq 0 \quad (i = 1, \ldots, s)$  \\
    \textbf{Result:} a decomposition $f = g_1 \cdot q_1 + \ldots + g_s \cdot q_s + r$ with $q_1, \ldots, q_s, r \in k[x_1, \dots, x_n]$ 
    
\begin{lstpseudo}
r:= 0      % remainder under construction
p := f     % to be divided away 
q_i := 0 (i = 1, ..., s)   % quotients under construction 
invariant of the while loop: f = g_1 * q_1 + ... + g_s * q_s + p + r
while p != 0: 
    if not QuotientUpdated():
        r := r + LT(p) 
        p := p - LT(p) 
    end if 
  end while
return (q_1, \ldots, q_s, r)  

QuotiensUpdated():  % tries to divide away the LT(p)
    for i = 1, ..., s: 
        if LT(p) is divisible by LT(g_i): 
            T:= LT(p)/LT(g_i) 
            q_i := q_i + T
            p := p - T * g_i
            return True
        end if 
    end for 
    return False
\end{lstpseudo}
\end{algorithm} 

\begin{example}[title = Example 2.2.1]
        $f 'x^2 \cdot y + x \cdot y^2 + y^2$ to be divided by $(g_1, g_2)$ with the respect to the lex order. 
        \begin{align*}
            g_1 & = x \cdot y -1 \\ 
            g_2 & = y^2 -1
        \end{align*}

    \paragraph{Solution:}
    \begin{alignat*}{2}
        f = &g_1 \cdot 0 && + g_2 \cdot 0 + \underbrace{(x^2 \cdot y + x \cdot y^2 + y^2)}_{\mathclap{\text{to be divided}}} + \underbrace{0}_{r} \\ 
        f = &g_1 \cdot x && + g_2 \cdot 0 + (x \cdot y^2 + x + y^2) + 0 \\
        f = &g_1 \cdot (x + y) && + g_2 \cdot 0 + (x + y + y^2) + 0 \\ 
        f = &g_1 \cdot (x + y) && + g_2 \cdot 0 + (y + y^2) + x \\
        f = &g_1 \cdot (x + y) && + g_2 \cdot y + (y + 1) + x \\
        f = &g_1 \cdot \underbrace{(x + y)}_{q_1} && + g_2 \cdot \underbrace{y}_{q_2} + 0 + \underbrace{(x + y + 1)}_{remainder}
    \end{alignat*}
\end{example}

\begin{example}[title = Example from School]
    \begin{alignat*}{2}
        1732 & = 3 \cdot 0 &&+ 1732 \\ 
             & = 3 \cdot 500 && + 232 \\ 
             & = 3 \cdot 570 && + 22 \\
             & = 3 \cdot 577 && + 1 \\ 
    \end{alignat*}
\end{example} 

\clearpage 

\begin{center}
  \textcolor{gray}{------------------------------------------- 10.11.26 --------------------------------------------} \\[0.3em]
\end{center}
\vspace{0.3em} 

\begin{theorem}[title = Theorem 2.2.2]
  Let $\succ$ be a monomial order on $\mathbb{Z}^n_{\geq 0}$ and let $G  = (g_1, \ldots, g_s)$ be a list of $s$ polynomials from $k[x_1, \dots, x_n]$. Then, every $f \in k[x_1, \dots, x_n]$ 
  admits a representation $f = g_1 \cdot  q_1 + \ldots + g_s \cdot q_s + r$ for some $q_1, \ldots, q_s, r \in k[x_1, \dots, x_n]$ such that for every $i = 1, \ldots, s$ the following hold:
  \begin{itemize}
    \item[a)] no term of $r$ is divisible by $LT(g_i)$ 
    \item[b)] $mdeg(f) \succeq mdeg(g_i \cdot q_i)$  
  \end{itemize} 

  \paragraph{Proof:}  The desired representation is obtained by the multivatiate long division of $f$ by $G$. We need to explain why this long division algorithm terminates. 

  \begin{minipage}{0.49\textwidth}
      The algorithm starts with $p=f$. 
      \[LT(T \cdot g_i) = T \cdot LT(g_i) = \frac{LT(p)}{LT(g_i)} \cdot LT(g_i) = LT(p)\] 
  \end{minipage}
  \hspace*{0.7cm}
  \begin{minipage}{0.39\textwidth}
\begin{lstpseudo}
if LT(p) is divisible by LT(g_i): 
    T := LT(p)/LT(g_i)
    q_i := q_i + T
    p := p - T * g_i
\end{lstpseudo}
  \end{minipage}

  The algorithm maintains the invariant $f = g_1 \cdot q_1 + \ldots + g_s \cdot q_s + p + r$ starting with $q_1 = \ldots = q_s = 0, \ p = f, \ r =0$. 

  In each iteration of the algorithm $mdeg \ p$ gets smaller, either due to a successful update of $q_1, \ldots, q_s,\ p$ or by transfering $LT(p)$ into the remainder. 
 
  As $\succ$ is a monomial order, every descending chain with the respect to this order is finite. So, the descrending chain of $LT(p)$ as $p$ changes throughout 
  the algorithm is finite. This means the algorithm will terminate, giving a desired decomposition of $f$. \\

  We still need to check, if (a) and (b) are true. 

  \begin{itemize}
    \item[(a)] is true, because only the terms in $p$ that do not divide any of the leading terms of $g_1, \ldots, g_s$ are transfered from $p$ to $r$. 
    \item[(b)] To verify, first observe that during the execution one has $mdeg \ f \succeq mdeg \ p$, because with $p = f$ and $mdeg \ p$ gets lower with
    each update of $p$.  \\[0.3em]
    
    \begin{center}
        \begin{tikzcd}[remember picture, row sep=0.5cm]
          |[alias=top]| g_i \cdot q_i + \dots \\
          |[alias=bottom]| g_i \cdot (q_i + T) + \dots
        \end{tikzcd}

        \begin{tikzpicture}[remember picture, overlay]
          \draw[kGreen, <->, shorten <=1pt, shorten >=1pt]
            let \p1=(top.east), \p2=(bottom.east), \n1={max(\x1,\x2)+4pt} in
            (\n1,\y1) to[bend left=60]
              node[midway, right=2pt, align=left, font=\scriptsize, text=kGreen]
                {update}
            (\n1,\y2);
        \end{tikzpicture}
    \end{center}  \ \\[0.3em]
    When a new term $T = \frac{LT(p)}{LT(g_i)}$ is added to $g_i$, we have \[mdeg(g_i \cdot T) = mdeg(g_i \cdot \frac{LT(p)}{LT(g_i)}) = mdeg(g_i) + mdeg(p) - mdeg(g_i) = mdeg(p)\]
    $\Longrightarrow \quad mdeg \ f \succeq mdeg \ p = mdeg(g, T) \quad \longrightarrow$ all terms added to $g_i \cdot q_i$ by updating $q_i$ to $q_i +T$ have the
    multidegree at most $mdef \ f$. 
  \end{itemize}
  \begin{flushright}
    $\Box $
  \end{flushright}
\end{theorem}

\begin{example}[title = Example 2.2.3]
    We divide $f = xy^2 - x$ by $G = (g_1, g_2)$ with $g_1 = xy + 1$ and $g_2 = y^2 -1$ with respect to the lex order. 
    \begin{align*} 
        xy^2 - x & = g_1 \cdot 0 + g_2 \cdot 0 + (xy^2 - x) + 0 \\
                & = g_1 \cdot y + g_2 \cdot 0 + (-y - x) + 0 \\ 
                & = g_1 \cdot y + g_2 \cdot 0 + (-y -x) 
    \end{align*}
    with $q_1 = y, q_2 = 0$ and $r = (-y-x)$. 

    On the other hand, $f \in \langle g_1, g_2 \rangle$, because $f = \ x \cdot g_2$  
\end{example} 
\begin{example}[title = Example 2.2.4 continued]

    What happend if we divide $f$ by $G' = (g_2, g_1)$ with respect to the lex order? 
    \begin{align*}
            xy^2 - x & = g_2 \cdot 0 + g_1 \cdot 0 + (xy^2 - x) + 0 \\ 
                    &= g_2 \cdot x + g_1 \cdot 0 + 0 + 0 \\
    \end{align*} 
    with $g_1 = x, g_2 = 0$ and $r = 0$.

\end{example}  

It looks like testing the ideal membership $f \in \langle g_1, \ldots, g_s \rangle$ doesn't work for an arbitrary choice of $(g_1, \ldots, g_s)$. 
We will figure out the choices that work. 

\begin{definition}[title = Definition 2.2.4]
   Let $Rem(f, G)$ be the \textit{remainder of the long division} of $f$ by $G = (g_1, \ldots, g_s)$ with $f, g_1, \ldots, g_s \in k[x_1, \dots, x_n]$. 
   Again, this depends on the monomial order.
\end{definition} 

